\section{Purpose}

\textbf{[STIPULATION]} This document develops the complete Lagrangian formulation
of gravitational dynamics in FFGFT: the intrinsic time field
$\mathcal T=1/\max(m,\omega)$ as the fundamental object, the field equation for
spherical and extended systems, the coupling of all matter fields to the time
field through conformal transformation, the modified Dirac equation, and the
complete total Lagrangian density. The document thus provides the framework
designated as ``gravitational dynamics'' in A140 and recorded there only as a
bridge result. What remains open after this document is explicitly named.

\section{The Intrinsic Time Field}

\textbf{[B]} The fundamental object of gravitational dynamics in FFGFT is not the
metric but the intrinsic time field:
\[
  \boxed{\;\mathcal T(x,t) = \frac{1}{\max(m,\omega)}\;}
\]
in natural units ($\hbar=c=1$). $[\mathcal T]=[E^{-1}]$. The field is reciprocal
to the local energy -- heavy objects have small $\mathcal T$ (rapidly oscillating),
light ones have large $\mathcal T$. The time-mass duality $\tilde T\cdot m=1$
(A020/A060) is the zero-field equation of this structure.

\textbf{[B]} \textbf{Three fundamental field geometries:}

\begin{center}
\renewcommand{\arraystretch}{1.3}
{\small
\begin{tabular}{ll}
\toprule
Geometry & Parameters \\
\midrule
Localised spherical &
  $\beta=2Gm/r$ (dimensionless), $\xi=2\sqrt G\cdot m$ (dimensionless),
  $\mathcal T(r)=(1/m_0)(1-\beta)$ \\
Localised non-spherical &
  $\beta_{ij}=r_{0ij}/r$ (tensor), $\xi_{ij}=2\sqrt G\cdot I_{ij}$ \\
Infinite homogeneous &
  $\nabla^2 m=4\pi G\rho_0 m+\Lambda_T m$, $\xi_\text{eff}=\xi/2$ \\
\bottomrule
\end{tabular}
}
\renewcommand{\arraystretch}{1.0}
\end{center}

\textbf{[K]} For quantum computing applications and atomic scales
($R/r\gg10^{12}$) only the localised spherical geometry is relevant; the others
are for theoretical completeness.

\section{The Field Equation}

\textbf{[B]} For a point-mass source $\rho=m\delta^3(\mathbf r)$ the complete
geometric solution is:
\[
  m(r) = m_0\Bigl(1+\frac{2Gm}{r}\Bigr) = m_0(1+\beta),
  \qquad
  \mathcal T(r) = \frac{1}{m_0}(1+\beta)^{-1}\approx\frac{1}{m_0}(1-\beta).
\]
The factor 2 in $r_0=2Gm$ follows from the relativistic field structure and
coincides exactly with the Schwarzschild radius. The field equation:
\[
  \nabla^2 m = 4\pi G\,\rho(x,t)\cdot m.
  \qquad [E^3 = E^{-2}\cdot E^4\cdot E = E^3\;]\checkmark
\]
For infinite systems with background consistency:
\[
  \nabla^2 m = 4\pi G\rho_0 m + \Lambda_T m,
  \qquad \Lambda_T = -4\pi G\rho_0.
\]
The $\Lambda_T$ term generates the cosmic shielding effect $\xi_\text{eff}=\xi/2$.

\textbf{[B]} \textbf{No new objects.} Neither $r_0=2Gm$ nor
$\Lambda_T=-4\pi G\rho_0$ are independent quantities beside the scale structure
from $\xi$. $r_0$ is the characteristic length \emph{of the respective mass} --
the same reciprocity as $L_0$, only with $m$ instead of the Planck mass as
reference. $\Lambda_T$ is the same field equation, computed from the homogeneous
background instead of the point source. Both are reference-point choices, not new
physics -- the same logic as the reference-point table in A262.

\section{The Universal Scale Parameter from Higgs Matching}

\textbf{[K]} The fundamental scale parameter $\xi$ is uniquely determined by the
Higgs physics:
\[
  \boxed{\;\xi = \frac{\lambda_h^2\,v^2}{16\pi^3\,m_h^2}
  \approx 1{.}33\times10^{-4}\;}
\]
with Higgs self-coupling $\lambda_h\approx0{.}13$, Higgs VEV $v\approx246\,\text{GeV}$,
Higgs mass $m_h\approx125\,\text{GeV}$. Dimensional verification:
$[\xi]=[\lambda_h^2 v^2]/[16\pi^3 m_h^2]=[E^2]/[E^2]=[1]$. The condition
$\beta_T=1$ in natural units ($\hbar=c=\alpha_\text{EM}=\beta_T=1$) fixes $\xi$
uniquely and eliminates all free parameters.

\section{The Lagrangian Formulation}

\textbf{[B]} The time-field Lagrangian density:
\[
  \mathcal L_\text{time}
  = \sqrt{-g}\,\Bigl[
    \tfrac12\,g^{\mu\nu}\partial_\mu\mathcal T\,\partial_\nu\mathcal T
    - V(\mathcal T)
  \Bigr].
\]
Dimensional consistency: $[\mathcal L_\text{time}]=[E^{-4}]([E^2]+[E^4])=[E^0]$.

\textbf{[B]} \textbf{Matter field coupling through conformal transformation.} All
matter fields couple to the time field via:
\[
  g_{\mu\nu}\;\to\;\Omega^2(\mathcal T)\,g_{\mu\nu},
  \qquad\Omega(\mathcal T)=\frac{\mathcal T_0}{\mathcal T}\quad[\text{dimensionless}].
\]

Scalar fields:
\[
  \mathcal L_\phi
  = \sqrt{-g}\,\Omega^4(\mathcal T)\,
    \Bigl(\tfrac12 g^{\mu\nu}\partial_\mu\phi\,\partial_\nu\phi
    - \tfrac12 m^2\phi^2\Bigr).
  \qquad [\mathcal L_\phi]=[E^{-4}][1][E^4]=[E^0]\checkmark
\]

Fermion fields:
\[
  \mathcal L_\psi
  = \sqrt{-g}\,\Omega^4(\mathcal T)\,
    \bigl(i\bar\psi\gamma^\mu\partial_\mu\psi - m\bar\psi\psi\bigr).
  \qquad [\mathcal L_\psi]=[E^{-4}][1][E^4]=[E^0]\checkmark
\]

Higgs time-field coupling:
\[
  \mathcal L_{\text{Higgs-T}}
  = |\mathcal D_{\text{Higgs-T}}|^2 - V(\mathcal T,\Phi),
  \qquad
  \mathcal D_{\text{Higgs-T}}
  = \mathcal T(\partial_\mu+igA_\mu)\Phi + \Phi\,\partial_\mu\mathcal T.
\]
This delivers $\mathcal T=1/(y\langle\Phi\rangle)$ -- the connection between
time field and Higgs vacuum expectation value.

\textbf{[B]} \textbf{Complete total Lagrangian density:}
\[
  \mathcal L_\text{total}
  = \mathcal L_\text{time}
  + \mathcal L_\text{gauge}
  + \mathcal L_\phi
  + \mathcal L_\psi
  + \mathcal L_{\text{Higgs-T}},
\]
with gauge term $\mathcal L_\text{gauge}=\sqrt{-g}\,(-\tfrac14 F_{\mu\nu}F^{\mu\nu})$.
Each summand is dimensionally $[E^0]$.

\section{The Modified Dirac Equation}

\textbf{[B]} In the FFGFT time-field framework the Dirac equation reads:
\[
  \bigl[i\gamma^\mu(\partial_\mu+\Gamma_\mu^{(T)})-m(x,t)\bigr]\psi=0,
\]
with the time-field connection:
\[
  \Gamma_\mu^{(T)}
  = \frac{1}{\mathcal T}\partial_\mu\mathcal T
  = -\frac{\partial_\mu m}{m^2}.
\]
This equation reduces in the limit $\mathcal T\to\text{const}$
($\Gamma_\mu^{(T)}\to0$) to the standard Dirac equation. The connection to the
Schrödinger equation (A075) remains via the polar decomposition of the vacuum
field.

\textbf{[B]} \textbf{Modified Schrödinger equation:}
\[
  i\,\mathcal T\,\frac{\partial\Psi}{\partial t}
  + i\,\Psi\,\Bigl[\frac{\partial\mathcal T}{\partial t}
  + \mathbf v\cdot\nabla\mathcal T\Bigr]
  = \hat H\,\Psi.
\]
In the limit $\partial_\mu\mathcal T\to0$ this reduces to the standard form
$i\hbar\,\partial_t\Psi=\hat H\Psi$ with $\mathcal T=1/m\leftrightarrow\hbar$.

\section{Dimensional Consistency Overview}

\textbf{[K]}
\begin{center}
\renewcommand{\arraystretch}{1.3}
{\small
\begin{tabular}{lccl}
\toprule
Equation & LHS & RHS & Status \\
\midrule
Time-field definition & $[E^{-1}]$ & $[1/\max(m,\omega)]=[E^{-1}]$ & $\checkmark$ \\
Field equation & $[E^3]$ & $[4\pi G\rho m]=[E^3]$ & $\checkmark$ \\
$\beta$ parameter & $[1]$ & $[2Gm/r]=[1]$ & $\checkmark$ \\
$\xi$ from Higgs & $[1]$ & $[\lambda_h^2 v^2/(16\pi^3 m_h^2)]=[1]$ & $\checkmark$ \\
Lagrangian density & $[E^0]$ & all terms & $\checkmark$ \\
\bottomrule
\end{tabular}
}
\renewcommand{\arraystretch}{1.0}
\end{center}

\section{Testable Predictions}

\textbf{[K]} The Lagrangian framework gives specific predictions:
\[
  a_e^{(T0)} = \frac\alpha{2\pi}\cdot\xi^2\cdot I_\text{loop}
  \approx 2{.}34\times10^{-10}
\]
(QED correction to the anomalous magnetic moment). Rotation curves with modified
potential: $v^2(r)=GM/r+\kappa r^2$. Energy-dependent time delay:
$\Delta t=(\xi/c)(1/E_1-1/E_2)(2Gm/r)$.

\section{Verification Script}

\begin{itemize}[leftmargin=2em,itemsep=2pt,topsep=2pt]
  \item Dimensional consistency of all terms numerically; limiting case Dirac
    $\to$ standard; $\xi$ from Higgs matching:
    \texttt{python/A\_Serie\_Skripte/a142\_gravitation\_lagrange.py}
\end{itemize}

\section{Symbol Glossary}

Symbols used throughout the entire A-series are explained in A010.
Specific to this document:

\begin{center}
\renewcommand{\arraystretch}{1.3}
{\small\begin{tabular}{lp{9cm}}
\toprule
Symbol & Meaning \\
\midrule
$\mathcal{T}(x,t)$ & Intrinsic time field (dynamic scalar field) \\
$\Omega$ & Conformal transformation factor \\
$\Gamma_\mu^{(T)}=\partial_\mu m/m^2$ & Time-field connection (connection coefficient) \\
$\mathcal{L}_\text{time}$ & Lagrangian density of the time field \\
$G_{\mu\nu}$ & Einstein tensor \\
$T_{\mu\nu}$ & Energy-momentum tensor \\
\bottomrule
\end{tabular}}
\renewcommand{\arraystretch}{1.0}
\end{center}

\section{Open Questions}

\textbf{A separate quantisation of gravitation is not required in FFGFT.}
$G$ is not an independent field quantity but follows intrinsically from $\xi$:
$G=\xi^2/(4m_\text{char})$. Gravitation is the projection of the time-field
geometry, not an independent quantum field. The question ``when is gravitation
quantised?'' does not arise in this framework -- it is already answered by the
basic structure. What remains: the translation of the classical field equations
into explicit loop corrections for scenarios beyond the Schwarzschild limit.

\textbf{The transition to measurable predictions beyond the Schwarzschild case is
not complete.} The QED correction (above) is computable; predictions for strong
fields, gravitational waves, or the post-Newtonian sector are not systematically
worked out.

\textbf{The connection to the projection chain (A090) is missing.} How the
gravitational sector fits into the type-I/II/III structure -- whether it is a
type-II loss or a type-I equivalent -- is not spelled out.

\section{Supersedes}

This document subsumes: Dok.~067 (time-field Lagrangian density with dimensional
consistency, matter-field coupling through conformal transformations, spherically
symmetric solutions, Higgs coupling, $\xi$ from Higgs matching, complete total
Lagrangian density, modified Dirac equation, experimental predictions). \emph{Not}
taken over: the section on a static universe from Dok.~067 -- it does not belong
in the A-series core (A010).

\section{Use of AI Tools}

This document was created with the assistance of AI tools. The questions,
stipulations, and all substantive decisions come from the author; the tools
formulate, compute, and create verification scripts.
Responsibility for content and errors rests entirely with the author.
Full wording of the rules in A010.
